\documentclass[17pt,a4paper,landscape]{extarticle}
\usepackage[margin=2.35cm,top=0.12cm,left=1.05cm]{geometry}
\usepackage{amsmath}
\usepackage{amssymb}
\usepackage{tasks}
\usepackage{tikz}
\usepackage{xcolor}
\usepackage[sfdefault,lf]{FiraSans}
\renewcommand{\familydefault}{\sfdefault}
\setlength{\parindent}{0pt}
\pagestyle{empty}
\settasks{label=\textbf{\arabic*.}, label-width=2.2em, item-indent=2.95em, column-sep=2.1cm, after-item-skip=4.7em}
\renewcommand{\arraystretch}{1.15}
\everymath{\displaystyle}

\begin{document}
\sffamily
\boldmath

\boldmath

{\LARGE \textbf{Using tables to explore patterns and graphs — Excellence}}\\[0.2em]
Use table patterns to justify rules, compare graphs, and predict unknown values.

\begin{tasks}(3)
\task The table is \mbox{$(0,4)$}, \mbox{$(1,7)$}, \mbox{$(2,10)$}, \mbox{$(3,13)$}. Write the rule, then explain how the table shows the graph has constant gradient.

\task The table is \mbox{$(1,12)$}, \mbox{$(2,9)$}, \mbox{$(3,6)$}, \mbox{$(4,3)$}. Write the rule and predict where the graph crosses the \mbox{$x$}-axis.

\task The table is \mbox{$(0,1)$}, \mbox{$(1,4)$}, \mbox{$(2,9)$}, \mbox{$(3,16)$}. Explain why these points would not lie on one straight line.

\task {Complete the table for \mbox{$y=3x-5$}. Then state both intercepts.\\[0.3em]
\begin{tabular}{|c|c|c|c|c|}
\hline
$x$ & $0$ & $1$ & $2$ & $5$ \\\hline
$y$ &  &  &  &  \\\hline
\end{tabular}}

\task Complete the table for \mbox{$y=18-2x$}. Then decide whether the graph reaches \mbox{$y=0$} before or after \mbox{$x=10$}.

\task A machine rule makes the table \mbox{$(0,2)$}, \mbox{$(2,6)$}, \mbox{$(4,10)$}, \mbox{$(6,14)$}. Write the rule in the form \mbox{$y=mx+c$}.

\task Two tables are shown. Table A: \mbox{$(0,1)$}, \mbox{$(1,3)$}, \mbox{$(2,5)$}. Table B: \mbox{$(0,1)$}, \mbox{$(1,4)$}, \mbox{$(2,7)$}. Which graph is steeper? Explain using the change in \mbox{$y$}.

\task A student says the table \mbox{$(0,7)$}, \mbox{$(1,5)$}, \mbox{$(2,3)$}, \mbox{$(3,1)$} has rule \mbox{$y=7-2x$}. Check the rule using the table, then find the point for \mbox{$x=5$}.

\task The table is \mbox{$(0,-1)$}, \mbox{$(1,2)$}, \mbox{$(2,5)$}, \mbox{$(3,8)$}. Predict the point where \mbox{$y=14$}. Show the pattern you used.

\task The points in a table are \mbox{$(0,10)$}, \mbox{$(5,8)$}, \mbox{$(10,6)$}, \mbox{$(15,4)$}. Describe the graph and find a rule that fits these values.

\task A graph is made from the rule \mbox{$y=\frac{3}{2}x+1$}. Write a table for \mbox{$x=0,2,4,6$}, then explain why these values are useful for plotting.

\task The table is \mbox{$(1,3)$}, \mbox{$(2,6)$}, \mbox{$(3,11)$}, \mbox{$(4,18)$}. The first differences are not constant. What does that suggest about the graph?

\task One pattern uses \mbox{$y=2x+1$}. Another uses \mbox{$y=2x+5$}. Explain how the tables would show the graphs are parallel.

\task The table for a line includes \mbox{$(2,9)$} and \mbox{$(5,15)$}. Find the gradient from the table, then write a rule for the line.
\end{tasks}

\end{document}
